\phantomsection
\chapter*{Аннотация}
\addcontentsline{toc}{chapter}{Аннотация}

Отчет по выпускной квалификационной работе содержит [ХХ] страниц текста, [Х] рисунков, [Х] таблиц, [Х] источников литературы и [Х] приложений.

\textbf{Ключевые слова:} АГРЕГАТОР ДАННЫХ, АВТОМАТИЗИРОВАННЫЙ СБОР ДАННЫХ, РАСПРЕДЕЛЕННЫЕ СИСТЕМЫ, МИКРОСЕРВИСНАЯ АРХИТЕКТУРА, ОТКАЗОУСТОЙЧИВОСТЬ,\\БОЛЬШИЕ ЯЗЫКОВЫЕ МОДЕЛИ, МНОГОКАНАЛЬНОЕ ОПОВЕЩЕНИЕ, ВЕБ-ПРИЛОЖЕНИЕ, TELEGRAM-БОТ.

Объектом исследования является процесс поиска, агрегации и обработки информации анонсов профильных мероприятий в сфере информационных технологий.

Предметом исследования выступает программная платформа для автоматизированного сбора, интеллектуального анализа и персонализированной доставки анонсов мероприятий.

Цель выпускной квалификационной работы — проектирование и разработка информационной системы, позволяющей снизить временные затраты специалистов на поиск релевантных событий в разрозненных источниках и избавить пользователей от информационной перегрузки.

В ходе выполнения работы была спроектирована распределенная отказоустойчивая архитектура системы. 
Разработана подсистема фонового сбора данных, способная автоматически извлекать текстовую информацию из открытых каналов в мессенджере Telegram и тематических Web-сайтов. 
Для суммаризации неформатированных текстов анонсов и выделения ключевых параметров событий (даты, места проведения, стоимости) успешно интегрированы методы обработки естественного языка на базе современных больших языковых моделей (YandexGPT).

Реализован алгоритм отбора анонсов событий по заданным критериям, и создан маршрутизатор для многоканальной рассылки сформированных подборок в мессенджере Telegram и по электронной почте. 
Спроектировано и разработано клиентское веб-приложение, предоставляющее удобный графический интерфейс для настройки профиля, управления источниками информации и просмотра ленты анонсов.

В результате разработана платформа, готовая к практическому применению. 
Практическая значимость работы заключается в создании инструмента, который автоматизирует рутинный процесс поиска информации, преобразует разрозненные тексты в единый стандартизированный формат и предоставляет специалистам удобный способ планирования профессионального досуга.

